% cap01/1.5_justificacion.tex — Justificación, delimitación y limitaciones.
\section{Justificación y Delimitación}
\subsection{Importancia}
La trascendencia y pertinencia de la presente investigación radican en proveer al ecosistema de voluntariado y a las Organizaciones No Gubernamentales (ONG) un instrumento de base tecnológica altamente escalable, estandarizado y seguro, que moderniza drásticamente su gestión. Al proporcionar una solución integral como SaaS, en la que la ONG no asume los exorbitantes costos de mantenimiento de servidores o de seguridad, se reducen de forma significativa las brechas tecnológicas del tercer sector en un país en vías de desarrollo, lo que apalanca la digitalización inclusiva e inserta a estas instituciones en la economía global del dato.

\subsection{Justificación Técnica y Tecnológica}
A nivel estrictamente ingenieril, el proyecto halla su fundamentación en la convergencia orquestada de paradigmas modernos de arquitectura de software y seguridad ofensiva/defensiva. El diseño contempla arquitecturas \textit{serverless}, la programación declarativa en interfaces (\textit{React Fiber}), así como estrategias profundas de aislamiento (\textit{Tenant Data Bleeding Prevention}). La centralización de la lógica en esquemas fuertemente tipados de base de datos, soportados bajo el principio ACID de PostgreSQL 16 y envueltos por políticas RLS y \textit{triggers} de \textit{Change Data Capture} (CDC), eleva la calidad arquitectónica. Adicionalmente, la adopción nativa de la metodología \textit{Design Documentation System} (DDS) para coordinar con \textit{Large Language Models} (LLMs) marca un precedente de vanguardia en la ingeniería reversa de requerimientos y asegura un \textit{Single Source of Truth} (SSOT).

\subsection{Justificación Social}
El tercer sector maneja la delicada labor de intervenir ante crisis humanitarias, brindar salud y proporcionar bienestar a comunidades excluidas. Sin embargo, su capacidad multiplicadora suele atrofiarse bajo la carga burocrática. Proveer a las organizaciones de una herramienta centralizada, segura y matemáticamente trazable fortalece radicalmente su resiliencia administrativa. Esto permite que dirijan un mayor volumen de sus recursos financieros y de esfuerzo humano directamente hacia sus misiones sociales en lugar de tareas operativas administrativas, maximizando su impacto filantrópico en el territorio nacional.

\subsection{Justificación Académica}
En el ámbito académico y en el corazón de la epistemología de la Ingeniería de Sistemas, el diseño topológico, la programación sistémica y la posterior validación experimental de "Democra" constituyen un insumo invaluable y un hito referencial. El trabajo compila la metodología \textit{end-to-end} (desde requisitos hasta operaciones \textit{cloud}), documenta matemáticamente las estrategias de seguridad perimetral de un SaaS genuino y elabora análisis estadísticos (alfa de Cronbach, cobertura de c\'odigo), posicionándose como un manual de ingeniería de altísimo rigor académico aplicable a arquitecturas en la nube y ciberseguridad defensiva.

\subsection{Delimitación}
\textbf{Espacial:} La traza de ejecución, pruebas y depuración del presente trabajo está centralizada e inmersa enteramente en el repositorio privado del código fuente del proyecto "Democra" y en su despliegue técnico en la plataforma de servicios en la nube (Supabase Cloud). \textbf{Temporal:} Corresponde estrictamente al desarrollo, revisión y estado arquitectónico del software tal cual es referenciado a lo largo del periodo de desarrollo comprendido en el año 2026. \textbf{Temática:} El análisis es interdisciplinar y cruza las fronteras cognitivas de la Ingeniería de Software Avanzada, Arquitectura de Componentes Web (SaaS y \textit{multi-tenant}), Seguridad de la Información e infraestructuras \textit{Zero-Trust} y Aseguramiento Cuantitativo de la Calidad del Software.

\subsection{Limitaciones y Alcance Restrictivo}
Conscientes del rigor epistemológico requerido para un trabajo de tesis de sistemas, se precisan las siguientes limitaciones teóricas y operativas del estudio empírico:
\begin{enumerate}
    \item \textbf{Limitante Metodológica (Enfoque DDS):} Dado el uso intrínseco del enfoque híbrido documental (DDS) que asiste al ciclo de ingeniería con IA, la recuperación y el análisis del modelo se priorizan sobre el estudio de los metadatos de los flujos operacionales. El análisis forense abarca la arquitectura de lectura del código y no necesariamente un \textit{post mortem} tras la salida pública del producto; es decir, no abarca una observación dinámica en el nicho comercial.
    \item \textbf{Validación Acotada (QA Sandbox):} La cuantificación técnica de la latencia, cobertura de código fuente y evaluaciones de métricas de experiencia mediante pruebas automatizadas han sido controladas paramétricamente en \textit{sandboxes} o escenarios de ensayo, y no han sido sometidas a una simulación de estrés DDoS externa (\textit{stress testing} con tráfico \textit{black hat} en producción).
    \item \textbf{Naturaleza Evolutiva del Ecosistema:} Como corolario de cualquier investigación basada en un \textit{stack} tecnológico en vías de hiperdesarrollo (\textit{TypeScript}, \textit{React Fiber 19}, \textit{Supabase PostgreSQL}), algunos métodos sintácticos, topologías o declaraciones expuestas poseen una dependencia temporal y están sujetas a deprecación si los mantenedores primarios liberan \textit{breaking changes} futuros a nivel global.
\end{enumerate}
